% Source: https://artofproblemsolving.com/community/c3h1383876_amc_10a_2017_15
% Stable identity: aops-topic-1383876-q15
\textbf{Solution.} With probability $\dfrac{1}{2}$ Laurent's number is in $[2017,4034]$, and then it is certainly at least Chlo\'e's number. With probability $\dfrac{1}{2}$ it is in $[0,2017]$, and then the two numbers are independent and uniform on the same interval, so by symmetry each is the larger with probability $\dfrac{1}{2}$ (ties have probability $0$). Hence
\[
P=\dfrac{1}{2}\cdot1+\dfrac{1}{2}\cdot\dfrac{1}{2}=\boxed{\dfrac{3}{4}},\quad\textbf{(C)}.
\]
Equivalently, the point (Chlo\'e, Laurent) is uniform on a $2017\times4034$ rectangle, and the region below the line $L=C$ is a right triangle of area $\dfrac{1}{2}\cdot2017^2$, one quarter of the rectangle's area.

\textbf{The error.} The numbers are \emph{real}, not integers: there are infinitely many choices, and a uniform choice from an interval cannot be replaced by a uniform choice of an integer. Even within the integer model there is a slip: when Chlo\'e chooses $2017$, Laurent must choose from $2018,\dots,4034$, which is $2017$ numbers, not $2018$, so the sum should be $2017+2018+\cdots+4034$ and the integer model gives
\[
\dfrac{2017+2018+\cdots+4034}{2018\cdot4035}=\dfrac{6051}{8070}\approx0.74981
\]
(the student's $\dfrac{3026\cdot2017}{2018\cdot4035}\approx0.74966$). Either way the result is close to but not equal to $\dfrac{3}{4}$: ties and endpoints matter for integers but have probability $0$ for real numbers. (If the integer model is refined to finer and finer grids, its answer tends to $\dfrac{3}{4}$.)
